Los experimentos corrieron en un equipo con procesador AMD Ryzen 5
4600H (6 núcleos, 12 hilos, 2.6 GHz), 16 GB de RAM DDR4, Windows 11, GPU NVIDIA GeForce GTX 1650, y se compiló en MSYS2 MINGW64.
El código se compiló con \texttt{g++} 16.1.0. El tiempo se midió con \texttt{chrono}; la memoria, con
el pico histórico del \textit{Working Set} de Windows
(\texttt{PeakWorkingSetSize}, vía \texttt{GetProcessMemoryInfo}), algo
específico de esta plataforma, así que si alguien intenta replicar
esto en Linux o macOS esa columna simplemente va a salir en cero.

Para ordenamiento se generaron arreglos de tamaño $n=10^1,10^3,10^5,10^7$,
en orden ascendente, descendente y aleatorio, con dominios $D_1=\{0,\dots,9\}$
y $D_7=\{0,\dots,10^7\}$, 3 muestras por combinación. Para matrices,
tamaños $n=2^4,2^6,2^8,2^{10}$, tipos densa/diagonal/dispersa, dominios
$D_0=\{0,1\}$ y $D_{10}=\{0,\dots,9\}$, también 3 muestras. Todo se
generó con semilla fija para poder repetir los experimentos. Además se
puso un límite de 3 minutos por caso, porque sin eso Quick Sort con
arreglos ya ordenados simplemente nunca termina (su peor caso
$O(n^2)$, que de hecho se activó varias veces en las mediciones).